% =============================================================================
%  Block:   INTEGRATOR -- system integration, budget, top-level testbench
%  Dir:     integrator/
%  Images:  integrator/images/  ->  \includegraphics{integrator/images/<file>}
%  Built from RF_PROJECT.tex via % =============================================================================
%  Block:   INTEGRATOR -- system integration, budget, top-level testbench
%  Dir:     integrator/
%  Images:  integrator/images/  ->  \includegraphics{integrator/images/<file>}
%  Built from RF_PROJECT.tex via % =============================================================================
%  Block:   INTEGRATOR -- system integration, budget, top-level testbench
%  Dir:     integrator/
%  Images:  integrator/images/  ->  \includegraphics{integrator/images/<file>}
%  Built from RF_PROJECT.tex via % =============================================================================
%  Block:   INTEGRATOR -- system integration, budget, top-level testbench
%  Dir:     integrator/
%  Images:  integrator/images/  ->  \includegraphics{integrator/images/<file>}
%  Built from RF_PROJECT.tex via \include{integrator/integrator}
% =============================================================================

\chapter{System Integration}
\label{ch:integrator}

\blockowner{Wojciech Miśkowicz}

% -----------------------------------------------------------------------------
\section{Schematic}
Block diagram of the RX and TX chains with the shared LO. Every block-to-block
connection is labelled with its interface ID I1--I10 (table in
section~\ref{sec:integrator-interfaces}).

\figplaceholder{TRX block diagram with interfaces I1--I10}{fig:integrator-block-diagram}

% -----------------------------------------------------------------------------
\section{Requirements}
Source: project specification and requirements changeover v0.2.

\begin{requirements}
	\req{SYS-01}{Technology}{TSMC 40\,nm, \texttt{tsmcN40}, Cadence + SpectreRF}{\specref{2--3}}
	\req{SYS-02}{Operating frequency}{6\,GHz, 20\,MHz channel}{\specref{2--3}}
	\req{SYS-03}{Architecture}{direct conversion, one LO shared by RX and TX (TDD)}{\specref{2--3}}
	\req{SYS-04}{Supply}{VDD = 1.1\,V for all blocks; VDD\_PA = 2\,V}{\specref{2--3}}
	\req{SYS-05}{RX input / TX output}{50\,$\Omega$}{\specref{2--3}}
	\req{SYS-06}{RX noise figure}{$< 10$\,dB}{\specref{2--3}}
	\req{SYS-07}{RX input signal range}{$-82$\,dBm (sensitivity)\dots$-30$\,dBm}{\specref{2--3}}
	\req{SYS-08}{TX output power}{$+16$\,dBm into 50\,$\Omega$, 8\,dB back-off from $P_\mathrm{1dB}$}{\specref{2--3}}
	\req{SYS-09}{Test signals}{one tone (gain, noise), two tones (IP3), QPSK/QAM (LNA and PA required)}{\specref{2--3}}
	\req{SYS-10}{Devices and bias}{real \texttt{tsmcN40} devices, bias inside each block; baluns/transformers may be ideal with a realistic $Q$}{\specref{2--3}}
	\req{SYS-11}{Coupling rule}{see the coupling rules below and the interface table}{\specref{2--3}}
\end{requirements}

\subsection{Derived system budget}\label{sec:integrator-budget-req}
These rows are not in the specification. They follow from SYS-06\dots SYS-08 and
from the converter models in section~\ref{sec:integrator-converters}. Every block
requirement table refers back to them. The derivations are in
section~\ref{sec:integrator-budget}.

\begin{requirements}
	\req{BUD-01}{RX noise figure (budget)}{NF$_\mathrm{SSB} \leq 7.5$\,dB at maximum gain (estimate 7.3\,dB, 2.7\,dB margin to SYS-06); measured SSB with one tone, single (I) mixer}{SYS-06}
	\req{BUD-02}{RX linearity}{$P_\mathrm{1dB,in} \geq -22$\,dBm and IIP3 $\geq -12$\,dBm at minimum gain, referred to the RX input (estimate $-11.4$\,dBm)}{SYS-07}
	\req{BUD-03}{RX voltage gain (pin $\to$ ADC)}{21\dots77\,dB (LNA 15 + DMIX 8 + VGA $-2$\dots54); nominal AGC target 100\,mVrms diff.\ at the ADC}{SYS-07}
	\req{BUD-04}{RX IIP2}{$\geq +35$\,dBm referred to the RX input}{SYS-03}
	\req{BUD-05}{TX chain $OP_\mathrm{1dB}$}{$\geq +24$\,dBm at the antenna; every stage before the PA compresses $\geq 2$\,dB later than the PA (referred to the PA output)}{SYS-08}
	\req{BUD-06}{TX nominal level plan}{DAC 70\,mVrms diff.\ $\to$ $+16$\,dBm at the antenna (section~\ref{tab:integrator-tx-levels})}{SYS-08}
	\req{BUD-07}{LO amplitude at each mixer}{1.0\,Vpp diff.\ nominal, 0.8\dots1.2\,Vpp diff.\ over PVT}{SYS-03}
	\req{BUD-08}{Power, RX mode}{$\leq 45$\,mW from 1.1\,V (LNA + DMIX + VGA + LO)}{\specref{3}}
	\req{BUD-09}{Power, TX mode}{$\leq 400$\,mW (UMIX + DRV + LO from 1.1\,V, PA from 2\,V)}{\specref{3}}
\end{requirements}

\subsection{Coupling rules}
\begin{enumerate}
	\item If the driving stage already has an inductive load, that load is made a
	      \textbf{transformer/balun}. It also gives AC coupling and single-ended to
	      differential conversion. The \textbf{driving} block owns it.
	\item \textbf{Bias across a transformer.} The secondary of an interface transformer
	      is \textbf{DC-floating}: its centre tap is only AC-grounded (MIM capacitor
	      $\geq 2$\,pF, owned by the driving block) and has no DC path to VDD or ground.
	      The \textbf{receiving} block sets the DC level of its own input gates through
	      its own bias resistors ($\geq 5$\,k$\Omega$ per side, so they do not load the
	      tank). No bias voltage crosses a block boundary (\specref{4}).
	\item Otherwise use a \textbf{MIM capacitor plus bias resistor} in the \textbf{receiving} block.
	\item No coil is added only for coupling.
	\item Baseband: DMIX$\to$VGA is AC-coupled (removes the DC offset from LO
	      self-mixing). The baseband edges (VGA$\to$ADC, DAC$\to$UMIX) are
	      DC-coupled with a common mode of \textbf{0.55\,V}.
	\item Impedances: 50\,$\Omega$ ports are given as $Z$ plus $|S_{11}|$/$|S_{22}|$;
	      on-chip ports as $R \parallel C$. The \textbf{receiving} block owns the number;
	      the driving block meets its own specification into that load.
\end{enumerate}

\subsection{Interfaces}\label{sec:integrator-interfaces}
The agreed numbers at each interface are the contract between blocks. The LO drives
I3 and I7 together, so the LO output load is the sum of both mixer inputs.

\par\medskip\noindent\small
\begin{tabularx}{\textwidth}{@{}l l l l L l@{}}
	\toprule
	\textbf{ID} & \textbf{From $\to$ To} & \textbf{Freq.} & \textbf{Type} & \textbf{Coupling} & \textbf{Owner} \\
	\midrule
	I1 & Antenna (50\,$\Omega$) $\to$ LNA & 6\,GHz & SE & AC: series C in the LNA input match & LNA \\
	I2 & LNA $\to$ DMIX & 6\,GHz & SE $\to$ diff. & AC: \textbf{transformer used as the LNA load}; secondary DC-floating, DMIX biases its own gates & LNA \\
	I3 & LO $\to$ DMIX & 6\,GHz & diff. & AC: C + R at the DMIX LO input & DMIX \\
	I4 & DMIX $\to$ VGA & BB $\leq 10$\,MHz & diff. & AC: C + R, $f_\mathrm{c} \leq 100$\,kHz & VGA \\
	I5 & VGA $\to$ ADC load & BB & diff. & \textbf{DC}, common mode 0.55\,V & -- \\
	I6 & DAC $\to$ UMIX & BB & diff. & \textbf{DC}, common mode 0.55\,V & -- \\
	I7 & LO $\to$ UMIX & 6\,GHz & diff. & AC: C + R at the UMIX LO input & UMIX \\
	I8 & UMIX $\to$ DRV & 6\,GHz & diff. & AC: \textbf{transformer used as the UMIX load} (= BPF tank, rule~1); secondary DC-floating. Proposed, UMIX and DRV owners confirm & UMIX \\
	I9 & DRV $\to$ PA & 6\,GHz & diff. & AC: \textbf{transformer used as the DRV load} (resonates the large PA $C_\mathrm{in}$); secondary DC-floating, PA biases its own gates & DRV \\
	I10 & PA $\to$ antenna (50\,$\Omega$) & 6\,GHz & diff. $\to$ SE & output transformer / matching network & PA \\
	\bottomrule
\end{tabularx}\par\normalsize

\medskip
\noindent Levels at each interface (CW test tone, nominal; ``diff.'' = differential,
``pk'' = peak, ``Vpp'' = peak-to-peak):

\par\medskip\noindent\small
\begin{tabularx}{\textwidth}{@{}l P{0.17\textwidth} P{0.25\textwidth} L@{}}
	\toprule
	\textbf{ID} & \textbf{DC level} & \textbf{Signal (min \dots max)} & \textbf{Load (owned by the receiver)} \\
	\midrule
	I1 & LNA gate $\approx 0.45$\,V (internal) & $-82$\dots$-30$\,dBm (50\,$\mu$Vpp\dots20\,mVpp at the pin) & 50\,$\Omega$, $|S_{11}| \leq -10$\,dB \\
	I2 & 0.55\,V, set by DMIX & 0.28\dots112\,mVpp diff. & $\geq 2$\,k$\Omega$ $\parallel$ $\leq 80$\,fF diff. \\
	I3, I7 & set by each mixer & 1.0\,Vpp diff.\ (0.8\dots1.2) & per mixer $\geq 5$\,k$\Omega$ $\parallel$ $\leq 80$\,fF diff. \\
	I4 & DMIX $\approx 0.7$\,V; VGA 0.55\,V & 0.71\dots283\,mVpp diff. & $\geq 50$\,k$\Omega$ $\parallel$ $\leq 1$\,pF diff.\ (after 20\,pF series C per side) \\
	I5 & \textbf{0.55\,V $\pm$ 30\,mV} & 283\,mVpp diff.\ (100\,mVrms) nominal, 1.0\,Vpp diff.\ full scale & ADC: 10\,k$\Omega$ diff.\ $\parallel$ 1\,pF per side \\
	I6 & \textbf{0.55\,V $\pm$ 50\,mV} & 200\,mVpp diff.\ (70\,mVrms) nominal, 1.0\,Vpp diff.\ full scale & $\geq 10$\,k$\Omega$ $\parallel$ $\leq 1$\,pF diff.; DAC $R_\mathrm{s} = 200\,\Omega$ diff. \\
	I8 & DRV gate $\approx 0.5$\,V, set by DRV & 0.2\,Vpp diff.\ (0.1\,V pk) & $\geq 1$\,k$\Omega$ $\parallel$ $\leq 150$\,fF diff. \\
	I9 & PA gate $\approx 0.45$\,V, set by PA & 1.0\,Vpp diff.\ (0.5\,V pk) & $\geq 100\,\Omega$ $\parallel$ $\leq 1.0$\,pF diff. \\
	I10 & drain at VDD\_PA = 2\,V & $+16$\,dBm (2.0\,V pk); $OP_\mathrm{1dB} \geq +24$\,dBm & 50\,$\Omega$, $|S_{22}| \leq -10$\,dB \\
	\bottomrule
\end{tabularx}\par\normalsize

% -----------------------------------------------------------------------------
\section{Overview}

\subsection{Top-level testbench}
\placeholder{\texttt{trx\_top} hierarchy, config views, stimulus and loads}

\figplaceholder{Top-level testbench \texttt{trx\_top}}{fig:integrator-testbench}

\subsection{Converter models at the baseband edges}\label{sec:integrator-converters}
The project has no ADC or DAC. The testbench models them as below. The values
are typical for a 20\,MHz Wi-Fi transceiver in a 1.1\,V, 40\,nm process.

\par\medskip\noindent\small
\begin{tabularx}{\textwidth}{@{}P{0.3\textwidth} L L@{}}
	\toprule
	\textbf{Parameter} & \textbf{ADC (RX, I5)} & \textbf{DAC (TX, I6)} \\
	\midrule
	Type & 10-bit SAR, 80\,MS/s (2$\times$ oversampled) & 10-bit current-steering, 160\,MS/s, with reconstruction filter \\
	Full scale & 1.0\,Vpp diff.\ (0.5\,V pk diff.) & 1.0\,Vpp diff. \\
	Common mode & 0.55\,V $\pm$ 30\,mV accepted & 0.55\,V $\pm$ 50\,mV delivered \\
	ENOB / SNR & 9 bit / 56\,dB $\to$ 89\,nV/$\sqrt{\mathrm{Hz}}$ diff.\ input-referred & 9 bit \\
	Nominal level & 100\,mVrms diff.\ ($-11$\,dBFS): OFDM PAPR 10\,dB + 4\,dB AGC margin & 70\,mVrms diff.\ ($-14$\,dBFS): set by UMIX linearity \\
	Model in \texttt{trx\_top} & 5\,k$\Omega$ per side to 0.55\,V + 1\,pF per side to ground ($R \approx 1/(f_\mathrm{s} C_\mathrm{s})$) & ideal source, 100\,$\Omega$ per side, CM 0.55\,V \\
	\bottomrule
\end{tabularx}\par\normalsize

\subsection{RX budget}\label{sec:integrator-budget}

\paragraph{Gain and level plan.}
Voltage gains are defined at the pins: LNA from the matched input pin (single-ended)
to its differential output; mixer conversion gain from RF diff.\ amplitude to BB diff.\
amplitude. The AGC target at the ADC is 100\,mVrms diff.\ (CW-equivalent).

\par\medskip\noindent\small
\begin{tabularx}{\textwidth}{@{}l c L L L@{}}
	\toprule
	\textbf{Node} & \textbf{Gain} & \textbf{$-82$\,dBm} & \textbf{$-30$\,dBm} & \textbf{$-22$\,dBm ($P_\mathrm{1dB}$ target)} \\
	\midrule
	RX input pin (SE) & -- & 50\,$\mu$Vpp & 20\,mVpp & 50\,mVpp \\
	LNA out / DMIX in (diff.) & 15\,dB & 0.28\,mVpp & 112\,mVpp & 283\,mVpp \\
	DMIX out / VGA in (diff.) & 8\,dB & 0.71\,mVpp & 283\,mVpp & 710\,mVpp \\
	VGA out / ADC (diff.) & $+52$ / $0$ / $0$\,dB & 283\,mVpp & 283\,mVpp & 710\,mVpp (FS 1.0\,Vpp) \\
	\bottomrule
\end{tabularx}\par\normalsize

\medskip
\noindent The required VGA range is 52\,dB (0\dots52\,dB). The VGA gives
$-2$\dots$+54$\,dB, so 2\,dB margin is left at each end for the PVT spread of the
LNA and DMIX gain. The LNA has one gain mode (decision D5), so the whole range is in
the VGA. At $-22$\,dBm with the gain still set for $-30$\,dBm the signal is 3\,dB
below the ADC full scale, so the ADC does not clip before the analog chain compresses.

\paragraph{Noise.}
All contributions are referred to the RX input pin as noise voltage density.
The source term is $kTR_\mathrm{s} = 2.0\cdot10^{-19}$\,V$^2$/Hz (0.447\,nV/$\sqrt{\mathrm{Hz}}$,
matched 50\,$\Omega$, $T_0 = 290$\,K). A block with input-referred density $v_n$
after voltage gain $G$ contributes $v_n^2/G^2$. NF is SSB (one mixer, the noise of
both sidebands folds onto the wanted one), which is 3\,dB worse than the DSB value
that an I/Q receiver would show. VGA at maximum gain:

\par\medskip\noindent\small
\begin{tabularx}{\textwidth}{@{}l l l L r@{}}
	\toprule
	\textbf{Block} & \textbf{Gain before} & \textbf{Own noise (budget)} & \textbf{Referred to input pin} & \textbf{Share} \\
	\midrule
	Source 50\,$\Omega$ & -- & -- & 0.447\,nV/$\sqrt{\mathrm{Hz}}$ & 19\,\% \\
	LNA & 0\,dB & NF $\leq 2$\,dB & 0.342\,nV/$\sqrt{\mathrm{Hz}}$ & 11\,\% \\
	DMIX & 15\,dB & $\leq 4.0$\,nV/$\sqrt{\mathrm{Hz}}$ diff.\ (SSB) & 0.711\,nV/$\sqrt{\mathrm{Hz}}$ & 47\,\% \\
	VGA + LPF & 23\,dB & $\leq 7.0$\,nV/$\sqrt{\mathrm{Hz}}$ diff. & 0.496\,nV/$\sqrt{\mathrm{Hz}}$ & 23\,\% \\
	ADC & 77\,dB & 89\,nV/$\sqrt{\mathrm{Hz}}$ diff. & 0.013\,nV/$\sqrt{\mathrm{Hz}}$ & $< 0.1$\,\% \\
	\midrule
	\textbf{Chain} & & & \textbf{NF$_\mathrm{SSB}$ = 7.3\,dB} (DSB $\approx 4.3$\,dB) & 100\,\% \\
	\bottomrule
\end{tabularx}\par\normalsize

\medskip
\noindent The input noise floor in 20\,MHz is $-93.7$\,dBm, so the SNR at $-82$\,dBm is
11.7\,dB. With a single-mode 15\,dB LNA the DMIX dominates (47\,\%), so its
noise budget is the tightest one in the RX. The 2.7\,dB margin to SYS-06 covers the
transformer loss at I2, real inductor $Q$ in the LNA and the PVT corners. At minimum VGA
gain the NF rises to $\approx 28$\,dB, but the input is then $\geq -30$\,dBm and the
SNR is still $\approx 43$\,dB.

\paragraph{Linearity.}
Two-tone IIP3 of each block is referred to the RX input and summed as
$1/\mathrm{IIP3} = \sum 1/\mathrm{IIP3}_{i,\mathrm{ref}}$ (in mW). The worst case is
the minimum VGA gain, where the full $-30$\,dBm signal reaches the VGA. Block IIP3 on
on-chip nodes is given as the \textbf{per-tone peak differential voltage at the block
input} at the intercept point.

\par\medskip\noindent\small
\begin{tabularx}{\textwidth}{@{}l L l r@{}}
	\toprule
	\textbf{Block} & \textbf{Requirement} & \textbf{Referred to RX input} & \textbf{Share} \\
	\midrule
	LNA & IIP3 $\geq -4$\,dBm (50\,$\Omega$) & $-4.0$\,dBm & 18\,\% \\
	DMIX & IIP3 $\geq 0.7$\,V pk diff.\ per tone at the RF input & $-8.1$\,dBm & 46\,\% \\
	VGA (code 0, $-2$\,dB) & OIP3 $\geq 1.6$\,V pk diff.\ per tone & $-6.9$\,dBm & 35\,\% \\
	\midrule
	\textbf{Chain} & & \textbf{$-11.4$\,dBm} & 100\,\% \\
	\bottomrule
\end{tabularx}\par\normalsize

\medskip
\noindent The requirement is IIP3 $\geq -12$\,dBm, from $P_\mathrm{1dB,in} \geq -22$\,dBm
($-30$\,dBm + 8\,dB) and IIP3 $\approx P_\mathrm{1dB} + 9.6$\,dB. The margin is only
0.6\,dB, so \textbf{linearity, not noise, is the critical RX budget}. If the DMIX
cannot reach 0.7\,V pk, the first remedy is a switchable LNA gain step
(e.g.\ 15/3\,dB, revisiting D5). It costs no noise at sensitivity and gains
$\approx 6$\,dB of IIP3.

\paragraph{IIP2 and DC offset.}
In direct conversion an AM interferer gives IM2 at DC and in band. BUD-04
($+35$\,dBm) keeps IM2 from a $-30$\,dBm two-tone signal 65\,dB below the signal. In
practice the DMIX mismatch sets IIP2. The DC offset from LO self-mixing is blocked
by the AC coupling at I4 ($f_\mathrm{c} \leq 100$\,kHz, below the first Wi-Fi
subcarrier at 312.5\,kHz). The VGA must cancel its own offset, because at 54\,dB a
2\,mV input offset would saturate the output.

\paragraph{Baseband filter.}
The LPF is a 3rd-order Butterworth with $f_{-3\mathrm{dB}} = 12$\,MHz. It costs
0.9\,dB at 9.4\,MHz (edge of the occupied band), gives 13.5\,dB at 20\,MHz (adjacent
channel) and 46\,dB at 70\,MHz (the first ADC alias at $f_\mathrm{s} - 10$\,MHz).

\subsection{TX budget}\label{tab:integrator-tx-levels}
The TX is specified in voltages up to the PA input and in dBm at the antenna. Each
stage before the PA has $\geq 10$\,dB back-off at the nominal level, so the PA alone
sets the chain $P_\mathrm{1dB}$. Without I/Q both sidebands ($f_\mathrm{LO} \pm f_\mathrm{BB}$)
leave the mixer at equal level. The $+16$\,dBm applies to the wanted sideband, and the
PA and DRV are also characterised alone with an RF source.

\par\medskip\noindent\small
\begin{tabularx}{\textwidth}{@{}l l L L@{}}
	\toprule
	\textbf{Node} & \textbf{Gain} & \textbf{Nominal (CW)} & \textbf{Compression requirement} \\
	\midrule
	DAC out / UMIX in & -- & 70\,mVrms = 0.1\,V pk diff. & FS 0.5\,V pk diff. \\
	UMIX + BPF out / DRV in & 0\,dB (wanted sideband) & 0.1\,V pk diff. & $OP_\mathrm{1dB} \geq 0.4$\,V pk diff.\ (12\,dB back-off) \\
	DRV out / PA in & 14\,dB & 0.5\,V pk diff. & $OP_\mathrm{1dB} \geq 1.6$\,V pk diff.\ (10\,dB back-off) \\
	PA out (50\,$\Omega$, SE) & 12\,dB (2.0\,V/0.5\,V) & $+16$\,dBm = 2.0\,V pk & $OP_\mathrm{1dB} \geq +24$\,dBm (5.0\,V pk); target $+25$\,dBm \\
	\bottomrule
\end{tabularx}\par\normalsize

\medskip
\noindent At the PA $P_\mathrm{1dB}$ its input is $0.5 \cdot 10^{9/20} = 1.41$\,V pk diff.
The DRV must deliver this with $\geq 1$\,dB to spare, so it needs 1.6\,V pk. The DAC sits
at $-14$\,dBFS. A 10-bit DAC still gives $\approx 42$\,dB SNR there, well above the
$\approx 25$\,dB needed for 64-QAM, and the UMIX input stays small enough for the
0.55\,V common mode at 1.1\,V supply.

\subsection{Power budget}

\par\medskip\noindent\small
\begin{tabularx}{\textwidth}{@{}l l r r L@{}}
	\toprule
	\textbf{Mode} & \textbf{Block} & \textbf{$I$ [mA]} & \textbf{$P$ [mW]} & \textbf{Basis} \\
	\midrule
	RX & LNA & 8 & 8.8 & NF 2\,dB at 6\,GHz needs $g_m \approx 50$--60\,mS \\
	RX & DMIX & 6 & 6.6 & Gilbert cell, $\approx 2.5$\,mA per side in the $g_m$ stage \\
	RX & VGA + LPF & 10 & 11 & 3 stages + LPF + offset cancellation \\
	RX & LO & 13 & 14.3 & LO-16/17 (owner's value) \\
	\textbf{RX} & \textbf{total} & \textbf{37} & \textbf{40.7} & BUD-08 $\leq 45$\,mW \\
	\midrule
	TX & UMIX & 6 & 6.6 & as DMIX \\
	TX & DRV & 25 & 27.5 & class A, 1.6\,V pk into $\geq 100\,\Omega$ diff.\ + transformer loss \\
	TX & PA & 175 (2\,V) & 350 & PAE $\geq 10$\,\% at $+16$\,dBm \\
	TX & LO & 13 & 14.3 & as RX \\
	\textbf{TX} & \textbf{total} & & \textbf{398} & BUD-09 $\leq 400$\,mW \\
	\bottomrule
\end{tabularx}\par\normalsize

\subsection{Budget and results}
\placeholder{parameter / budget / stub / design / status table}

% -----------------------------------------------------------------------------
\section{Changelog}

\begin{changelog}
	\change{2026-11-12}{integrator}{Requirements v0.1 (budget and interfaces) from design proposal}{Starting point for meeting~2}{all}{proposed}
	\change{2026-11-17}{integrator}{Requirements v0.1 $\to$ v0.2 (standard table)}{Changeover, see requirements\_changeover.md}{all}{proposed}
	\change{2026-11-27}{integrator}{Requirements v0.4: derived budget BUD-01\dots09, ADC/DAC models, interface levels, RX noise/linearity cascade, TX level plan, power budget}{Fill TBD values so that each block can design alone}{all}{proposed}
	\change{2026-11-27}{integrator}{Transformer centre tap: ``biases the receiver'' $\to$ secondary DC-floating, receiver biases its own gates}{No bias voltage may cross a block boundary (\specref{4})}{LNA, DMIX, UMIX, DRV, PA}{proposed}
	\change{2026-11-27}{integrator}{I8: transformer or C + R $\to$ transformer as UMIX load (= BPF tank)}{Rule~1: UMIX already has an inductive load}{UMIX, DRV}{proposed}
	\change{2026-11-27}{integrator}{Supply current and power: TBD $\to$ budget values (D10 reversed)}{Power budget is a team deliverable (\specref{3}, item~5)}{all}{proposed}
\end{changelog}

% =============================================================================

\chapter{System Integration}
\label{ch:integrator}

\blockowner{Wojciech Miśkowicz}

% -----------------------------------------------------------------------------
\section{Schematic}
Block diagram of the RX and TX chains with the shared LO. Every block-to-block
connection is labelled with its interface ID I1--I10 (table in
section~\ref{sec:integrator-interfaces}).

\figplaceholder{TRX block diagram with interfaces I1--I10}{fig:integrator-block-diagram}

% -----------------------------------------------------------------------------
\section{Requirements}
Source: project specification and requirements changeover v0.2.

\begin{requirements}
	\req{SYS-01}{Technology}{TSMC 40\,nm, \texttt{tsmcN40}, Cadence + SpectreRF}{\specref{2--3}}
	\req{SYS-02}{Operating frequency}{6\,GHz, 20\,MHz channel}{\specref{2--3}}
	\req{SYS-03}{Architecture}{direct conversion, one LO shared by RX and TX (TDD)}{\specref{2--3}}
	\req{SYS-04}{Supply}{VDD = 1.1\,V for all blocks; VDD\_PA = 2\,V}{\specref{2--3}}
	\req{SYS-05}{RX input / TX output}{50\,$\Omega$}{\specref{2--3}}
	\req{SYS-06}{RX noise figure}{$< 10$\,dB}{\specref{2--3}}
	\req{SYS-07}{RX input signal range}{$-82$\,dBm (sensitivity)\dots$-30$\,dBm}{\specref{2--3}}
	\req{SYS-08}{TX output power}{$+16$\,dBm into 50\,$\Omega$, 8\,dB back-off from $P_\mathrm{1dB}$}{\specref{2--3}}
	\req{SYS-09}{Test signals}{one tone (gain, noise), two tones (IP3), QPSK/QAM (LNA and PA required)}{\specref{2--3}}
	\req{SYS-10}{Devices and bias}{real \texttt{tsmcN40} devices, bias inside each block; baluns/transformers may be ideal with a realistic $Q$}{\specref{2--3}}
	\req{SYS-11}{Coupling rule}{see the coupling rules below and the interface table}{\specref{2--3}}
\end{requirements}

\subsection{Derived system budget}\label{sec:integrator-budget-req}
These rows are not in the specification. They follow from SYS-06\dots SYS-08 and
from the converter models in section~\ref{sec:integrator-converters}. Every block
requirement table refers back to them. The derivations are in
section~\ref{sec:integrator-budget}.

\begin{requirements}
	\req{BUD-01}{RX noise figure (budget)}{NF$_\mathrm{SSB} \leq 7.5$\,dB at maximum gain (estimate 7.3\,dB, 2.7\,dB margin to SYS-06); measured SSB with one tone, single (I) mixer}{SYS-06}
	\req{BUD-02}{RX linearity}{$P_\mathrm{1dB,in} \geq -22$\,dBm and IIP3 $\geq -12$\,dBm at minimum gain, referred to the RX input (estimate $-11.4$\,dBm)}{SYS-07}
	\req{BUD-03}{RX voltage gain (pin $\to$ ADC)}{21\dots77\,dB (LNA 15 + DMIX 8 + VGA $-2$\dots54); nominal AGC target 100\,mVrms diff.\ at the ADC}{SYS-07}
	\req{BUD-04}{RX IIP2}{$\geq +35$\,dBm referred to the RX input}{SYS-03}
	\req{BUD-05}{TX chain $OP_\mathrm{1dB}$}{$\geq +24$\,dBm at the antenna; every stage before the PA compresses $\geq 2$\,dB later than the PA (referred to the PA output)}{SYS-08}
	\req{BUD-06}{TX nominal level plan}{DAC 70\,mVrms diff.\ $\to$ $+16$\,dBm at the antenna (section~\ref{tab:integrator-tx-levels})}{SYS-08}
	\req{BUD-07}{LO amplitude at each mixer}{1.0\,Vpp diff.\ nominal, 0.8\dots1.2\,Vpp diff.\ over PVT}{SYS-03}
	\req{BUD-08}{Power, RX mode}{$\leq 45$\,mW from 1.1\,V (LNA + DMIX + VGA + LO)}{\specref{3}}
	\req{BUD-09}{Power, TX mode}{$\leq 400$\,mW (UMIX + DRV + LO from 1.1\,V, PA from 2\,V)}{\specref{3}}
\end{requirements}

\subsection{Coupling rules}
\begin{enumerate}
	\item If the driving stage already has an inductive load, that load is made a
	      \textbf{transformer/balun}. It also gives AC coupling and single-ended to
	      differential conversion. The \textbf{driving} block owns it.
	\item \textbf{Bias across a transformer.} The secondary of an interface transformer
	      is \textbf{DC-floating}: its centre tap is only AC-grounded (MIM capacitor
	      $\geq 2$\,pF, owned by the driving block) and has no DC path to VDD or ground.
	      The \textbf{receiving} block sets the DC level of its own input gates through
	      its own bias resistors ($\geq 5$\,k$\Omega$ per side, so they do not load the
	      tank). No bias voltage crosses a block boundary (\specref{4}).
	\item Otherwise use a \textbf{MIM capacitor plus bias resistor} in the \textbf{receiving} block.
	\item No coil is added only for coupling.
	\item Baseband: DMIX$\to$VGA is AC-coupled (removes the DC offset from LO
	      self-mixing). The baseband edges (VGA$\to$ADC, DAC$\to$UMIX) are
	      DC-coupled with a common mode of \textbf{0.55\,V}.
	\item Impedances: 50\,$\Omega$ ports are given as $Z$ plus $|S_{11}|$/$|S_{22}|$;
	      on-chip ports as $R \parallel C$. The \textbf{receiving} block owns the number;
	      the driving block meets its own specification into that load.
\end{enumerate}

\subsection{Interfaces}\label{sec:integrator-interfaces}
The agreed numbers at each interface are the contract between blocks. The LO drives
I3 and I7 together, so the LO output load is the sum of both mixer inputs.

\par\medskip\noindent\small
\begin{tabularx}{\textwidth}{@{}l l l l L l@{}}
	\toprule
	\textbf{ID} & \textbf{From $\to$ To} & \textbf{Freq.} & \textbf{Type} & \textbf{Coupling} & \textbf{Owner} \\
	\midrule
	I1 & Antenna (50\,$\Omega$) $\to$ LNA & 6\,GHz & SE & AC: series C in the LNA input match & LNA \\
	I2 & LNA $\to$ DMIX & 6\,GHz & SE $\to$ diff. & AC: \textbf{transformer used as the LNA load}; secondary DC-floating, DMIX biases its own gates & LNA \\
	I3 & LO $\to$ DMIX & 6\,GHz & diff. & AC: C + R at the DMIX LO input & DMIX \\
	I4 & DMIX $\to$ VGA & BB $\leq 10$\,MHz & diff. & AC: C + R, $f_\mathrm{c} \leq 100$\,kHz & VGA \\
	I5 & VGA $\to$ ADC load & BB & diff. & \textbf{DC}, common mode 0.55\,V & -- \\
	I6 & DAC $\to$ UMIX & BB & diff. & \textbf{DC}, common mode 0.55\,V & -- \\
	I7 & LO $\to$ UMIX & 6\,GHz & diff. & AC: C + R at the UMIX LO input & UMIX \\
	I8 & UMIX $\to$ DRV & 6\,GHz & diff. & AC: \textbf{transformer used as the UMIX load} (= BPF tank, rule~1); secondary DC-floating. Proposed, UMIX and DRV owners confirm & UMIX \\
	I9 & DRV $\to$ PA & 6\,GHz & diff. & AC: \textbf{transformer used as the DRV load} (resonates the large PA $C_\mathrm{in}$); secondary DC-floating, PA biases its own gates & DRV \\
	I10 & PA $\to$ antenna (50\,$\Omega$) & 6\,GHz & diff. $\to$ SE & output transformer / matching network & PA \\
	\bottomrule
\end{tabularx}\par\normalsize

\medskip
\noindent Levels at each interface (CW test tone, nominal; ``diff.'' = differential,
``pk'' = peak, ``Vpp'' = peak-to-peak):

\par\medskip\noindent\small
\begin{tabularx}{\textwidth}{@{}l P{0.17\textwidth} P{0.25\textwidth} L@{}}
	\toprule
	\textbf{ID} & \textbf{DC level} & \textbf{Signal (min \dots max)} & \textbf{Load (owned by the receiver)} \\
	\midrule
	I1 & LNA gate $\approx 0.45$\,V (internal) & $-82$\dots$-30$\,dBm (50\,$\mu$Vpp\dots20\,mVpp at the pin) & 50\,$\Omega$, $|S_{11}| \leq -10$\,dB \\
	I2 & 0.55\,V, set by DMIX & 0.28\dots112\,mVpp diff. & $\geq 2$\,k$\Omega$ $\parallel$ $\leq 80$\,fF diff. \\
	I3, I7 & set by each mixer & 1.0\,Vpp diff.\ (0.8\dots1.2) & per mixer $\geq 5$\,k$\Omega$ $\parallel$ $\leq 80$\,fF diff. \\
	I4 & DMIX $\approx 0.7$\,V; VGA 0.55\,V & 0.71\dots283\,mVpp diff. & $\geq 50$\,k$\Omega$ $\parallel$ $\leq 1$\,pF diff.\ (after 20\,pF series C per side) \\
	I5 & \textbf{0.55\,V $\pm$ 30\,mV} & 283\,mVpp diff.\ (100\,mVrms) nominal, 1.0\,Vpp diff.\ full scale & ADC: 10\,k$\Omega$ diff.\ $\parallel$ 1\,pF per side \\
	I6 & \textbf{0.55\,V $\pm$ 50\,mV} & 200\,mVpp diff.\ (70\,mVrms) nominal, 1.0\,Vpp diff.\ full scale & $\geq 10$\,k$\Omega$ $\parallel$ $\leq 1$\,pF diff.; DAC $R_\mathrm{s} = 200\,\Omega$ diff. \\
	I8 & DRV gate $\approx 0.5$\,V, set by DRV & 0.2\,Vpp diff.\ (0.1\,V pk) & $\geq 1$\,k$\Omega$ $\parallel$ $\leq 150$\,fF diff. \\
	I9 & PA gate $\approx 0.45$\,V, set by PA & 1.0\,Vpp diff.\ (0.5\,V pk) & $\geq 100\,\Omega$ $\parallel$ $\leq 1.0$\,pF diff. \\
	I10 & drain at VDD\_PA = 2\,V & $+16$\,dBm (2.0\,V pk); $OP_\mathrm{1dB} \geq +24$\,dBm & 50\,$\Omega$, $|S_{22}| \leq -10$\,dB \\
	\bottomrule
\end{tabularx}\par\normalsize

% -----------------------------------------------------------------------------
\section{Overview}

\subsection{Top-level testbench}
\placeholder{\texttt{trx\_top} hierarchy, config views, stimulus and loads}

\figplaceholder{Top-level testbench \texttt{trx\_top}}{fig:integrator-testbench}

\subsection{Converter models at the baseband edges}\label{sec:integrator-converters}
The project has no ADC or DAC. The testbench models them as below. The values
are typical for a 20\,MHz Wi-Fi transceiver in a 1.1\,V, 40\,nm process.

\par\medskip\noindent\small
\begin{tabularx}{\textwidth}{@{}P{0.3\textwidth} L L@{}}
	\toprule
	\textbf{Parameter} & \textbf{ADC (RX, I5)} & \textbf{DAC (TX, I6)} \\
	\midrule
	Type & 10-bit SAR, 80\,MS/s (2$\times$ oversampled) & 10-bit current-steering, 160\,MS/s, with reconstruction filter \\
	Full scale & 1.0\,Vpp diff.\ (0.5\,V pk diff.) & 1.0\,Vpp diff. \\
	Common mode & 0.55\,V $\pm$ 30\,mV accepted & 0.55\,V $\pm$ 50\,mV delivered \\
	ENOB / SNR & 9 bit / 56\,dB $\to$ 89\,nV/$\sqrt{\mathrm{Hz}}$ diff.\ input-referred & 9 bit \\
	Nominal level & 100\,mVrms diff.\ ($-11$\,dBFS): OFDM PAPR 10\,dB + 4\,dB AGC margin & 70\,mVrms diff.\ ($-14$\,dBFS): set by UMIX linearity \\
	Model in \texttt{trx\_top} & 5\,k$\Omega$ per side to 0.55\,V + 1\,pF per side to ground ($R \approx 1/(f_\mathrm{s} C_\mathrm{s})$) & ideal source, 100\,$\Omega$ per side, CM 0.55\,V \\
	\bottomrule
\end{tabularx}\par\normalsize

\subsection{RX budget}\label{sec:integrator-budget}

\paragraph{Gain and level plan.}
Voltage gains are defined at the pins: LNA from the matched input pin (single-ended)
to its differential output; mixer conversion gain from RF diff.\ amplitude to BB diff.\
amplitude. The AGC target at the ADC is 100\,mVrms diff.\ (CW-equivalent).

\par\medskip\noindent\small
\begin{tabularx}{\textwidth}{@{}l c L L L@{}}
	\toprule
	\textbf{Node} & \textbf{Gain} & \textbf{$-82$\,dBm} & \textbf{$-30$\,dBm} & \textbf{$-22$\,dBm ($P_\mathrm{1dB}$ target)} \\
	\midrule
	RX input pin (SE) & -- & 50\,$\mu$Vpp & 20\,mVpp & 50\,mVpp \\
	LNA out / DMIX in (diff.) & 15\,dB & 0.28\,mVpp & 112\,mVpp & 283\,mVpp \\
	DMIX out / VGA in (diff.) & 8\,dB & 0.71\,mVpp & 283\,mVpp & 710\,mVpp \\
	VGA out / ADC (diff.) & $+52$ / $0$ / $0$\,dB & 283\,mVpp & 283\,mVpp & 710\,mVpp (FS 1.0\,Vpp) \\
	\bottomrule
\end{tabularx}\par\normalsize

\medskip
\noindent The required VGA range is 52\,dB (0\dots52\,dB). The VGA gives
$-2$\dots$+54$\,dB, so 2\,dB margin is left at each end for the PVT spread of the
LNA and DMIX gain. The LNA has one gain mode (decision D5), so the whole range is in
the VGA. At $-22$\,dBm with the gain still set for $-30$\,dBm the signal is 3\,dB
below the ADC full scale, so the ADC does not clip before the analog chain compresses.

\paragraph{Noise.}
All contributions are referred to the RX input pin as noise voltage density.
The source term is $kTR_\mathrm{s} = 2.0\cdot10^{-19}$\,V$^2$/Hz (0.447\,nV/$\sqrt{\mathrm{Hz}}$,
matched 50\,$\Omega$, $T_0 = 290$\,K). A block with input-referred density $v_n$
after voltage gain $G$ contributes $v_n^2/G^2$. NF is SSB (one mixer, the noise of
both sidebands folds onto the wanted one), which is 3\,dB worse than the DSB value
that an I/Q receiver would show. VGA at maximum gain:

\par\medskip\noindent\small
\begin{tabularx}{\textwidth}{@{}l l l L r@{}}
	\toprule
	\textbf{Block} & \textbf{Gain before} & \textbf{Own noise (budget)} & \textbf{Referred to input pin} & \textbf{Share} \\
	\midrule
	Source 50\,$\Omega$ & -- & -- & 0.447\,nV/$\sqrt{\mathrm{Hz}}$ & 19\,\% \\
	LNA & 0\,dB & NF $\leq 2$\,dB & 0.342\,nV/$\sqrt{\mathrm{Hz}}$ & 11\,\% \\
	DMIX & 15\,dB & $\leq 4.0$\,nV/$\sqrt{\mathrm{Hz}}$ diff.\ (SSB) & 0.711\,nV/$\sqrt{\mathrm{Hz}}$ & 47\,\% \\
	VGA + LPF & 23\,dB & $\leq 7.0$\,nV/$\sqrt{\mathrm{Hz}}$ diff. & 0.496\,nV/$\sqrt{\mathrm{Hz}}$ & 23\,\% \\
	ADC & 77\,dB & 89\,nV/$\sqrt{\mathrm{Hz}}$ diff. & 0.013\,nV/$\sqrt{\mathrm{Hz}}$ & $< 0.1$\,\% \\
	\midrule
	\textbf{Chain} & & & \textbf{NF$_\mathrm{SSB}$ = 7.3\,dB} (DSB $\approx 4.3$\,dB) & 100\,\% \\
	\bottomrule
\end{tabularx}\par\normalsize

\medskip
\noindent The input noise floor in 20\,MHz is $-93.7$\,dBm, so the SNR at $-82$\,dBm is
11.7\,dB. With a single-mode 15\,dB LNA the DMIX dominates (47\,\%), so its
noise budget is the tightest one in the RX. The 2.7\,dB margin to SYS-06 covers the
transformer loss at I2, real inductor $Q$ in the LNA and the PVT corners. At minimum VGA
gain the NF rises to $\approx 28$\,dB, but the input is then $\geq -30$\,dBm and the
SNR is still $\approx 43$\,dB.

\paragraph{Linearity.}
Two-tone IIP3 of each block is referred to the RX input and summed as
$1/\mathrm{IIP3} = \sum 1/\mathrm{IIP3}_{i,\mathrm{ref}}$ (in mW). The worst case is
the minimum VGA gain, where the full $-30$\,dBm signal reaches the VGA. Block IIP3 on
on-chip nodes is given as the \textbf{per-tone peak differential voltage at the block
input} at the intercept point.

\par\medskip\noindent\small
\begin{tabularx}{\textwidth}{@{}l L l r@{}}
	\toprule
	\textbf{Block} & \textbf{Requirement} & \textbf{Referred to RX input} & \textbf{Share} \\
	\midrule
	LNA & IIP3 $\geq -4$\,dBm (50\,$\Omega$) & $-4.0$\,dBm & 18\,\% \\
	DMIX & IIP3 $\geq 0.7$\,V pk diff.\ per tone at the RF input & $-8.1$\,dBm & 46\,\% \\
	VGA (code 0, $-2$\,dB) & OIP3 $\geq 1.6$\,V pk diff.\ per tone & $-6.9$\,dBm & 35\,\% \\
	\midrule
	\textbf{Chain} & & \textbf{$-11.4$\,dBm} & 100\,\% \\
	\bottomrule
\end{tabularx}\par\normalsize

\medskip
\noindent The requirement is IIP3 $\geq -12$\,dBm, from $P_\mathrm{1dB,in} \geq -22$\,dBm
($-30$\,dBm + 8\,dB) and IIP3 $\approx P_\mathrm{1dB} + 9.6$\,dB. The margin is only
0.6\,dB, so \textbf{linearity, not noise, is the critical RX budget}. If the DMIX
cannot reach 0.7\,V pk, the first remedy is a switchable LNA gain step
(e.g.\ 15/3\,dB, revisiting D5). It costs no noise at sensitivity and gains
$\approx 6$\,dB of IIP3.

\paragraph{IIP2 and DC offset.}
In direct conversion an AM interferer gives IM2 at DC and in band. BUD-04
($+35$\,dBm) keeps IM2 from a $-30$\,dBm two-tone signal 65\,dB below the signal. In
practice the DMIX mismatch sets IIP2. The DC offset from LO self-mixing is blocked
by the AC coupling at I4 ($f_\mathrm{c} \leq 100$\,kHz, below the first Wi-Fi
subcarrier at 312.5\,kHz). The VGA must cancel its own offset, because at 54\,dB a
2\,mV input offset would saturate the output.

\paragraph{Baseband filter.}
The LPF is a 3rd-order Butterworth with $f_{-3\mathrm{dB}} = 12$\,MHz. It costs
0.9\,dB at 9.4\,MHz (edge of the occupied band), gives 13.5\,dB at 20\,MHz (adjacent
channel) and 46\,dB at 70\,MHz (the first ADC alias at $f_\mathrm{s} - 10$\,MHz).

\subsection{TX budget}\label{tab:integrator-tx-levels}
The TX is specified in voltages up to the PA input and in dBm at the antenna. Each
stage before the PA has $\geq 10$\,dB back-off at the nominal level, so the PA alone
sets the chain $P_\mathrm{1dB}$. Without I/Q both sidebands ($f_\mathrm{LO} \pm f_\mathrm{BB}$)
leave the mixer at equal level. The $+16$\,dBm applies to the wanted sideband, and the
PA and DRV are also characterised alone with an RF source.

\par\medskip\noindent\small
\begin{tabularx}{\textwidth}{@{}l l L L@{}}
	\toprule
	\textbf{Node} & \textbf{Gain} & \textbf{Nominal (CW)} & \textbf{Compression requirement} \\
	\midrule
	DAC out / UMIX in & -- & 70\,mVrms = 0.1\,V pk diff. & FS 0.5\,V pk diff. \\
	UMIX + BPF out / DRV in & 0\,dB (wanted sideband) & 0.1\,V pk diff. & $OP_\mathrm{1dB} \geq 0.4$\,V pk diff.\ (12\,dB back-off) \\
	DRV out / PA in & 14\,dB & 0.5\,V pk diff. & $OP_\mathrm{1dB} \geq 1.6$\,V pk diff.\ (10\,dB back-off) \\
	PA out (50\,$\Omega$, SE) & 12\,dB (2.0\,V/0.5\,V) & $+16$\,dBm = 2.0\,V pk & $OP_\mathrm{1dB} \geq +24$\,dBm (5.0\,V pk); target $+25$\,dBm \\
	\bottomrule
\end{tabularx}\par\normalsize

\medskip
\noindent At the PA $P_\mathrm{1dB}$ its input is $0.5 \cdot 10^{9/20} = 1.41$\,V pk diff.
The DRV must deliver this with $\geq 1$\,dB to spare, so it needs 1.6\,V pk. The DAC sits
at $-14$\,dBFS. A 10-bit DAC still gives $\approx 42$\,dB SNR there, well above the
$\approx 25$\,dB needed for 64-QAM, and the UMIX input stays small enough for the
0.55\,V common mode at 1.1\,V supply.

\subsection{Power budget}

\par\medskip\noindent\small
\begin{tabularx}{\textwidth}{@{}l l r r L@{}}
	\toprule
	\textbf{Mode} & \textbf{Block} & \textbf{$I$ [mA]} & \textbf{$P$ [mW]} & \textbf{Basis} \\
	\midrule
	RX & LNA & 8 & 8.8 & NF 2\,dB at 6\,GHz needs $g_m \approx 50$--60\,mS \\
	RX & DMIX & 6 & 6.6 & Gilbert cell, $\approx 2.5$\,mA per side in the $g_m$ stage \\
	RX & VGA + LPF & 10 & 11 & 3 stages + LPF + offset cancellation \\
	RX & LO & 13 & 14.3 & LO-16/17 (owner's value) \\
	\textbf{RX} & \textbf{total} & \textbf{37} & \textbf{40.7} & BUD-08 $\leq 45$\,mW \\
	\midrule
	TX & UMIX & 6 & 6.6 & as DMIX \\
	TX & DRV & 25 & 27.5 & class A, 1.6\,V pk into $\geq 100\,\Omega$ diff.\ + transformer loss \\
	TX & PA & 175 (2\,V) & 350 & PAE $\geq 10$\,\% at $+16$\,dBm \\
	TX & LO & 13 & 14.3 & as RX \\
	\textbf{TX} & \textbf{total} & & \textbf{398} & BUD-09 $\leq 400$\,mW \\
	\bottomrule
\end{tabularx}\par\normalsize

\subsection{Budget and results}
\placeholder{parameter / budget / stub / design / status table}

% -----------------------------------------------------------------------------
\section{Changelog}

\begin{changelog}
	\change{2026-11-12}{integrator}{Requirements v0.1 (budget and interfaces) from design proposal}{Starting point for meeting~2}{all}{proposed}
	\change{2026-11-17}{integrator}{Requirements v0.1 $\to$ v0.2 (standard table)}{Changeover, see requirements\_changeover.md}{all}{proposed}
	\change{2026-11-27}{integrator}{Requirements v0.4: derived budget BUD-01\dots09, ADC/DAC models, interface levels, RX noise/linearity cascade, TX level plan, power budget}{Fill TBD values so that each block can design alone}{all}{proposed}
	\change{2026-11-27}{integrator}{Transformer centre tap: ``biases the receiver'' $\to$ secondary DC-floating, receiver biases its own gates}{No bias voltage may cross a block boundary (\specref{4})}{LNA, DMIX, UMIX, DRV, PA}{proposed}
	\change{2026-11-27}{integrator}{I8: transformer or C + R $\to$ transformer as UMIX load (= BPF tank)}{Rule~1: UMIX already has an inductive load}{UMIX, DRV}{proposed}
	\change{2026-11-27}{integrator}{Supply current and power: TBD $\to$ budget values (D10 reversed)}{Power budget is a team deliverable (\specref{3}, item~5)}{all}{proposed}
\end{changelog}

% =============================================================================

\chapter{System Integration}
\label{ch:integrator}

\blockowner{Wojciech Miśkowicz}

% -----------------------------------------------------------------------------
\section{Schematic}
Block diagram of the RX and TX chains with the shared LO. Every block-to-block
connection is labelled with its interface ID I1--I10 (table in
section~\ref{sec:integrator-interfaces}).

\figplaceholder{TRX block diagram with interfaces I1--I10}{fig:integrator-block-diagram}

% -----------------------------------------------------------------------------
\section{Requirements}
Source: project specification and requirements changeover v0.2.

\begin{requirements}
	\req{SYS-01}{Technology}{TSMC 40\,nm, \texttt{tsmcN40}, Cadence + SpectreRF}{\specref{2--3}}
	\req{SYS-02}{Operating frequency}{6\,GHz, 20\,MHz channel}{\specref{2--3}}
	\req{SYS-03}{Architecture}{direct conversion, one LO shared by RX and TX (TDD)}{\specref{2--3}}
	\req{SYS-04}{Supply}{VDD = 1.1\,V for all blocks; VDD\_PA = 2\,V}{\specref{2--3}}
	\req{SYS-05}{RX input / TX output}{50\,$\Omega$}{\specref{2--3}}
	\req{SYS-06}{RX noise figure}{$< 10$\,dB}{\specref{2--3}}
	\req{SYS-07}{RX input signal range}{$-82$\,dBm (sensitivity)\dots$-30$\,dBm}{\specref{2--3}}
	\req{SYS-08}{TX output power}{$+16$\,dBm into 50\,$\Omega$, 8\,dB back-off from $P_\mathrm{1dB}$}{\specref{2--3}}
	\req{SYS-09}{Test signals}{one tone (gain, noise), two tones (IP3), QPSK/QAM (LNA and PA required)}{\specref{2--3}}
	\req{SYS-10}{Devices and bias}{real \texttt{tsmcN40} devices, bias inside each block; baluns/transformers may be ideal with a realistic $Q$}{\specref{2--3}}
	\req{SYS-11}{Coupling rule}{see the coupling rules below and the interface table}{\specref{2--3}}
\end{requirements}

\subsection{Derived system budget}\label{sec:integrator-budget-req}
These rows are not in the specification. They follow from SYS-06\dots SYS-08 and
from the converter models in section~\ref{sec:integrator-converters}. Every block
requirement table refers back to them. The derivations are in
section~\ref{sec:integrator-budget}.

\begin{requirements}
	\req{BUD-01}{RX noise figure (budget)}{NF$_\mathrm{SSB} \leq 7.5$\,dB at maximum gain (estimate 7.3\,dB, 2.7\,dB margin to SYS-06); measured SSB with one tone, single (I) mixer}{SYS-06}
	\req{BUD-02}{RX linearity}{$P_\mathrm{1dB,in} \geq -22$\,dBm and IIP3 $\geq -12$\,dBm at minimum gain, referred to the RX input (estimate $-11.4$\,dBm)}{SYS-07}
	\req{BUD-03}{RX voltage gain (pin $\to$ ADC)}{21\dots77\,dB (LNA 15 + DMIX 8 + VGA $-2$\dots54); nominal AGC target 100\,mVrms diff.\ at the ADC}{SYS-07}
	\req{BUD-04}{RX IIP2}{$\geq +35$\,dBm referred to the RX input}{SYS-03}
	\req{BUD-05}{TX chain $OP_\mathrm{1dB}$}{$\geq +24$\,dBm at the antenna; every stage before the PA compresses $\geq 2$\,dB later than the PA (referred to the PA output)}{SYS-08}
	\req{BUD-06}{TX nominal level plan}{DAC 70\,mVrms diff.\ $\to$ $+16$\,dBm at the antenna (section~\ref{tab:integrator-tx-levels})}{SYS-08}
	\req{BUD-07}{LO amplitude at each mixer}{1.0\,Vpp diff.\ nominal, 0.8\dots1.2\,Vpp diff.\ over PVT}{SYS-03}
	\req{BUD-08}{Power, RX mode}{$\leq 45$\,mW from 1.1\,V (LNA + DMIX + VGA + LO)}{\specref{3}}
	\req{BUD-09}{Power, TX mode}{$\leq 400$\,mW (UMIX + DRV + LO from 1.1\,V, PA from 2\,V)}{\specref{3}}
\end{requirements}

\subsection{Coupling rules}
\begin{enumerate}
	\item If the driving stage already has an inductive load, that load is made a
	      \textbf{transformer/balun}. It also gives AC coupling and single-ended to
	      differential conversion. The \textbf{driving} block owns it.
	\item \textbf{Bias across a transformer.} The secondary of an interface transformer
	      is \textbf{DC-floating}: its centre tap is only AC-grounded (MIM capacitor
	      $\geq 2$\,pF, owned by the driving block) and has no DC path to VDD or ground.
	      The \textbf{receiving} block sets the DC level of its own input gates through
	      its own bias resistors ($\geq 5$\,k$\Omega$ per side, so they do not load the
	      tank). No bias voltage crosses a block boundary (\specref{4}).
	\item Otherwise use a \textbf{MIM capacitor plus bias resistor} in the \textbf{receiving} block.
	\item No coil is added only for coupling.
	\item Baseband: DMIX$\to$VGA is AC-coupled (removes the DC offset from LO
	      self-mixing). The baseband edges (VGA$\to$ADC, DAC$\to$UMIX) are
	      DC-coupled with a common mode of \textbf{0.55\,V}.
	\item Impedances: 50\,$\Omega$ ports are given as $Z$ plus $|S_{11}|$/$|S_{22}|$;
	      on-chip ports as $R \parallel C$. The \textbf{receiving} block owns the number;
	      the driving block meets its own specification into that load.
\end{enumerate}

\subsection{Interfaces}\label{sec:integrator-interfaces}
The agreed numbers at each interface are the contract between blocks. The LO drives
I3 and I7 together, so the LO output load is the sum of both mixer inputs.

\par\medskip\noindent\small
\begin{tabularx}{\textwidth}{@{}l l l l L l@{}}
	\toprule
	\textbf{ID} & \textbf{From $\to$ To} & \textbf{Freq.} & \textbf{Type} & \textbf{Coupling} & \textbf{Owner} \\
	\midrule
	I1 & Antenna (50\,$\Omega$) $\to$ LNA & 6\,GHz & SE & AC: series C in the LNA input match & LNA \\
	I2 & LNA $\to$ DMIX & 6\,GHz & SE $\to$ diff. & AC: \textbf{transformer used as the LNA load}; secondary DC-floating, DMIX biases its own gates & LNA \\
	I3 & LO $\to$ DMIX & 6\,GHz & diff. & AC: C + R at the DMIX LO input & DMIX \\
	I4 & DMIX $\to$ VGA & BB $\leq 10$\,MHz & diff. & AC: C + R, $f_\mathrm{c} \leq 100$\,kHz & VGA \\
	I5 & VGA $\to$ ADC load & BB & diff. & \textbf{DC}, common mode 0.55\,V & -- \\
	I6 & DAC $\to$ UMIX & BB & diff. & \textbf{DC}, common mode 0.55\,V & -- \\
	I7 & LO $\to$ UMIX & 6\,GHz & diff. & AC: C + R at the UMIX LO input & UMIX \\
	I8 & UMIX $\to$ DRV & 6\,GHz & diff. & AC: \textbf{transformer used as the UMIX load} (= BPF tank, rule~1); secondary DC-floating. Proposed, UMIX and DRV owners confirm & UMIX \\
	I9 & DRV $\to$ PA & 6\,GHz & diff. & AC: \textbf{transformer used as the DRV load} (resonates the large PA $C_\mathrm{in}$); secondary DC-floating, PA biases its own gates & DRV \\
	I10 & PA $\to$ antenna (50\,$\Omega$) & 6\,GHz & diff. $\to$ SE & output transformer / matching network & PA \\
	\bottomrule
\end{tabularx}\par\normalsize

\medskip
\noindent Levels at each interface (CW test tone, nominal; ``diff.'' = differential,
``pk'' = peak, ``Vpp'' = peak-to-peak):

\par\medskip\noindent\small
\begin{tabularx}{\textwidth}{@{}l P{0.17\textwidth} P{0.25\textwidth} L@{}}
	\toprule
	\textbf{ID} & \textbf{DC level} & \textbf{Signal (min \dots max)} & \textbf{Load (owned by the receiver)} \\
	\midrule
	I1 & LNA gate $\approx 0.45$\,V (internal) & $-82$\dots$-30$\,dBm (50\,$\mu$Vpp\dots20\,mVpp at the pin) & 50\,$\Omega$, $|S_{11}| \leq -10$\,dB \\
	I2 & 0.55\,V, set by DMIX & 0.28\dots112\,mVpp diff. & $\geq 2$\,k$\Omega$ $\parallel$ $\leq 80$\,fF diff. \\
	I3, I7 & set by each mixer & 1.0\,Vpp diff.\ (0.8\dots1.2) & per mixer $\geq 5$\,k$\Omega$ $\parallel$ $\leq 80$\,fF diff. \\
	I4 & DMIX $\approx 0.7$\,V; VGA 0.55\,V & 0.71\dots283\,mVpp diff. & $\geq 50$\,k$\Omega$ $\parallel$ $\leq 1$\,pF diff.\ (after 20\,pF series C per side) \\
	I5 & \textbf{0.55\,V $\pm$ 30\,mV} & 283\,mVpp diff.\ (100\,mVrms) nominal, 1.0\,Vpp diff.\ full scale & ADC: 10\,k$\Omega$ diff.\ $\parallel$ 1\,pF per side \\
	I6 & \textbf{0.55\,V $\pm$ 50\,mV} & 200\,mVpp diff.\ (70\,mVrms) nominal, 1.0\,Vpp diff.\ full scale & $\geq 10$\,k$\Omega$ $\parallel$ $\leq 1$\,pF diff.; DAC $R_\mathrm{s} = 200\,\Omega$ diff. \\
	I8 & DRV gate $\approx 0.5$\,V, set by DRV & 0.2\,Vpp diff.\ (0.1\,V pk) & $\geq 1$\,k$\Omega$ $\parallel$ $\leq 150$\,fF diff. \\
	I9 & PA gate $\approx 0.45$\,V, set by PA & 1.0\,Vpp diff.\ (0.5\,V pk) & $\geq 100\,\Omega$ $\parallel$ $\leq 1.0$\,pF diff. \\
	I10 & drain at VDD\_PA = 2\,V & $+16$\,dBm (2.0\,V pk); $OP_\mathrm{1dB} \geq +24$\,dBm & 50\,$\Omega$, $|S_{22}| \leq -10$\,dB \\
	\bottomrule
\end{tabularx}\par\normalsize

% -----------------------------------------------------------------------------
\section{Overview}

\subsection{Top-level testbench}
\placeholder{\texttt{trx\_top} hierarchy, config views, stimulus and loads}

\figplaceholder{Top-level testbench \texttt{trx\_top}}{fig:integrator-testbench}

\subsection{Converter models at the baseband edges}\label{sec:integrator-converters}
The project has no ADC or DAC. The testbench models them as below. The values
are typical for a 20\,MHz Wi-Fi transceiver in a 1.1\,V, 40\,nm process.

\par\medskip\noindent\small
\begin{tabularx}{\textwidth}{@{}P{0.3\textwidth} L L@{}}
	\toprule
	\textbf{Parameter} & \textbf{ADC (RX, I5)} & \textbf{DAC (TX, I6)} \\
	\midrule
	Type & 10-bit SAR, 80\,MS/s (2$\times$ oversampled) & 10-bit current-steering, 160\,MS/s, with reconstruction filter \\
	Full scale & 1.0\,Vpp diff.\ (0.5\,V pk diff.) & 1.0\,Vpp diff. \\
	Common mode & 0.55\,V $\pm$ 30\,mV accepted & 0.55\,V $\pm$ 50\,mV delivered \\
	ENOB / SNR & 9 bit / 56\,dB $\to$ 89\,nV/$\sqrt{\mathrm{Hz}}$ diff.\ input-referred & 9 bit \\
	Nominal level & 100\,mVrms diff.\ ($-11$\,dBFS): OFDM PAPR 10\,dB + 4\,dB AGC margin & 70\,mVrms diff.\ ($-14$\,dBFS): set by UMIX linearity \\
	Model in \texttt{trx\_top} & 5\,k$\Omega$ per side to 0.55\,V + 1\,pF per side to ground ($R \approx 1/(f_\mathrm{s} C_\mathrm{s})$) & ideal source, 100\,$\Omega$ per side, CM 0.55\,V \\
	\bottomrule
\end{tabularx}\par\normalsize

\subsection{RX budget}\label{sec:integrator-budget}

\paragraph{Gain and level plan.}
Voltage gains are defined at the pins: LNA from the matched input pin (single-ended)
to its differential output; mixer conversion gain from RF diff.\ amplitude to BB diff.\
amplitude. The AGC target at the ADC is 100\,mVrms diff.\ (CW-equivalent).

\par\medskip\noindent\small
\begin{tabularx}{\textwidth}{@{}l c L L L@{}}
	\toprule
	\textbf{Node} & \textbf{Gain} & \textbf{$-82$\,dBm} & \textbf{$-30$\,dBm} & \textbf{$-22$\,dBm ($P_\mathrm{1dB}$ target)} \\
	\midrule
	RX input pin (SE) & -- & 50\,$\mu$Vpp & 20\,mVpp & 50\,mVpp \\
	LNA out / DMIX in (diff.) & 15\,dB & 0.28\,mVpp & 112\,mVpp & 283\,mVpp \\
	DMIX out / VGA in (diff.) & 8\,dB & 0.71\,mVpp & 283\,mVpp & 710\,mVpp \\
	VGA out / ADC (diff.) & $+52$ / $0$ / $0$\,dB & 283\,mVpp & 283\,mVpp & 710\,mVpp (FS 1.0\,Vpp) \\
	\bottomrule
\end{tabularx}\par\normalsize

\medskip
\noindent The required VGA range is 52\,dB (0\dots52\,dB). The VGA gives
$-2$\dots$+54$\,dB, so 2\,dB margin is left at each end for the PVT spread of the
LNA and DMIX gain. The LNA has one gain mode (decision D5), so the whole range is in
the VGA. At $-22$\,dBm with the gain still set for $-30$\,dBm the signal is 3\,dB
below the ADC full scale, so the ADC does not clip before the analog chain compresses.

\paragraph{Noise.}
All contributions are referred to the RX input pin as noise voltage density.
The source term is $kTR_\mathrm{s} = 2.0\cdot10^{-19}$\,V$^2$/Hz (0.447\,nV/$\sqrt{\mathrm{Hz}}$,
matched 50\,$\Omega$, $T_0 = 290$\,K). A block with input-referred density $v_n$
after voltage gain $G$ contributes $v_n^2/G^2$. NF is SSB (one mixer, the noise of
both sidebands folds onto the wanted one), which is 3\,dB worse than the DSB value
that an I/Q receiver would show. VGA at maximum gain:

\par\medskip\noindent\small
\begin{tabularx}{\textwidth}{@{}l l l L r@{}}
	\toprule
	\textbf{Block} & \textbf{Gain before} & \textbf{Own noise (budget)} & \textbf{Referred to input pin} & \textbf{Share} \\
	\midrule
	Source 50\,$\Omega$ & -- & -- & 0.447\,nV/$\sqrt{\mathrm{Hz}}$ & 19\,\% \\
	LNA & 0\,dB & NF $\leq 2$\,dB & 0.342\,nV/$\sqrt{\mathrm{Hz}}$ & 11\,\% \\
	DMIX & 15\,dB & $\leq 4.0$\,nV/$\sqrt{\mathrm{Hz}}$ diff.\ (SSB) & 0.711\,nV/$\sqrt{\mathrm{Hz}}$ & 47\,\% \\
	VGA + LPF & 23\,dB & $\leq 7.0$\,nV/$\sqrt{\mathrm{Hz}}$ diff. & 0.496\,nV/$\sqrt{\mathrm{Hz}}$ & 23\,\% \\
	ADC & 77\,dB & 89\,nV/$\sqrt{\mathrm{Hz}}$ diff. & 0.013\,nV/$\sqrt{\mathrm{Hz}}$ & $< 0.1$\,\% \\
	\midrule
	\textbf{Chain} & & & \textbf{NF$_\mathrm{SSB}$ = 7.3\,dB} (DSB $\approx 4.3$\,dB) & 100\,\% \\
	\bottomrule
\end{tabularx}\par\normalsize

\medskip
\noindent The input noise floor in 20\,MHz is $-93.7$\,dBm, so the SNR at $-82$\,dBm is
11.7\,dB. With a single-mode 15\,dB LNA the DMIX dominates (47\,\%), so its
noise budget is the tightest one in the RX. The 2.7\,dB margin to SYS-06 covers the
transformer loss at I2, real inductor $Q$ in the LNA and the PVT corners. At minimum VGA
gain the NF rises to $\approx 28$\,dB, but the input is then $\geq -30$\,dBm and the
SNR is still $\approx 43$\,dB.

\paragraph{Linearity.}
Two-tone IIP3 of each block is referred to the RX input and summed as
$1/\mathrm{IIP3} = \sum 1/\mathrm{IIP3}_{i,\mathrm{ref}}$ (in mW). The worst case is
the minimum VGA gain, where the full $-30$\,dBm signal reaches the VGA. Block IIP3 on
on-chip nodes is given as the \textbf{per-tone peak differential voltage at the block
input} at the intercept point.

\par\medskip\noindent\small
\begin{tabularx}{\textwidth}{@{}l L l r@{}}
	\toprule
	\textbf{Block} & \textbf{Requirement} & \textbf{Referred to RX input} & \textbf{Share} \\
	\midrule
	LNA & IIP3 $\geq -4$\,dBm (50\,$\Omega$) & $-4.0$\,dBm & 18\,\% \\
	DMIX & IIP3 $\geq 0.7$\,V pk diff.\ per tone at the RF input & $-8.1$\,dBm & 46\,\% \\
	VGA (code 0, $-2$\,dB) & OIP3 $\geq 1.6$\,V pk diff.\ per tone & $-6.9$\,dBm & 35\,\% \\
	\midrule
	\textbf{Chain} & & \textbf{$-11.4$\,dBm} & 100\,\% \\
	\bottomrule
\end{tabularx}\par\normalsize

\medskip
\noindent The requirement is IIP3 $\geq -12$\,dBm, from $P_\mathrm{1dB,in} \geq -22$\,dBm
($-30$\,dBm + 8\,dB) and IIP3 $\approx P_\mathrm{1dB} + 9.6$\,dB. The margin is only
0.6\,dB, so \textbf{linearity, not noise, is the critical RX budget}. If the DMIX
cannot reach 0.7\,V pk, the first remedy is a switchable LNA gain step
(e.g.\ 15/3\,dB, revisiting D5). It costs no noise at sensitivity and gains
$\approx 6$\,dB of IIP3.

\paragraph{IIP2 and DC offset.}
In direct conversion an AM interferer gives IM2 at DC and in band. BUD-04
($+35$\,dBm) keeps IM2 from a $-30$\,dBm two-tone signal 65\,dB below the signal. In
practice the DMIX mismatch sets IIP2. The DC offset from LO self-mixing is blocked
by the AC coupling at I4 ($f_\mathrm{c} \leq 100$\,kHz, below the first Wi-Fi
subcarrier at 312.5\,kHz). The VGA must cancel its own offset, because at 54\,dB a
2\,mV input offset would saturate the output.

\paragraph{Baseband filter.}
The LPF is a 3rd-order Butterworth with $f_{-3\mathrm{dB}} = 12$\,MHz. It costs
0.9\,dB at 9.4\,MHz (edge of the occupied band), gives 13.5\,dB at 20\,MHz (adjacent
channel) and 46\,dB at 70\,MHz (the first ADC alias at $f_\mathrm{s} - 10$\,MHz).

\subsection{TX budget}\label{tab:integrator-tx-levels}
The TX is specified in voltages up to the PA input and in dBm at the antenna. Each
stage before the PA has $\geq 10$\,dB back-off at the nominal level, so the PA alone
sets the chain $P_\mathrm{1dB}$. Without I/Q both sidebands ($f_\mathrm{LO} \pm f_\mathrm{BB}$)
leave the mixer at equal level. The $+16$\,dBm applies to the wanted sideband, and the
PA and DRV are also characterised alone with an RF source.

\par\medskip\noindent\small
\begin{tabularx}{\textwidth}{@{}l l L L@{}}
	\toprule
	\textbf{Node} & \textbf{Gain} & \textbf{Nominal (CW)} & \textbf{Compression requirement} \\
	\midrule
	DAC out / UMIX in & -- & 70\,mVrms = 0.1\,V pk diff. & FS 0.5\,V pk diff. \\
	UMIX + BPF out / DRV in & 0\,dB (wanted sideband) & 0.1\,V pk diff. & $OP_\mathrm{1dB} \geq 0.4$\,V pk diff.\ (12\,dB back-off) \\
	DRV out / PA in & 14\,dB & 0.5\,V pk diff. & $OP_\mathrm{1dB} \geq 1.6$\,V pk diff.\ (10\,dB back-off) \\
	PA out (50\,$\Omega$, SE) & 12\,dB (2.0\,V/0.5\,V) & $+16$\,dBm = 2.0\,V pk & $OP_\mathrm{1dB} \geq +24$\,dBm (5.0\,V pk); target $+25$\,dBm \\
	\bottomrule
\end{tabularx}\par\normalsize

\medskip
\noindent At the PA $P_\mathrm{1dB}$ its input is $0.5 \cdot 10^{9/20} = 1.41$\,V pk diff.
The DRV must deliver this with $\geq 1$\,dB to spare, so it needs 1.6\,V pk. The DAC sits
at $-14$\,dBFS. A 10-bit DAC still gives $\approx 42$\,dB SNR there, well above the
$\approx 25$\,dB needed for 64-QAM, and the UMIX input stays small enough for the
0.55\,V common mode at 1.1\,V supply.

\subsection{Power budget}

\par\medskip\noindent\small
\begin{tabularx}{\textwidth}{@{}l l r r L@{}}
	\toprule
	\textbf{Mode} & \textbf{Block} & \textbf{$I$ [mA]} & \textbf{$P$ [mW]} & \textbf{Basis} \\
	\midrule
	RX & LNA & 8 & 8.8 & NF 2\,dB at 6\,GHz needs $g_m \approx 50$--60\,mS \\
	RX & DMIX & 6 & 6.6 & Gilbert cell, $\approx 2.5$\,mA per side in the $g_m$ stage \\
	RX & VGA + LPF & 10 & 11 & 3 stages + LPF + offset cancellation \\
	RX & LO & 13 & 14.3 & LO-16/17 (owner's value) \\
	\textbf{RX} & \textbf{total} & \textbf{37} & \textbf{40.7} & BUD-08 $\leq 45$\,mW \\
	\midrule
	TX & UMIX & 6 & 6.6 & as DMIX \\
	TX & DRV & 25 & 27.5 & class A, 1.6\,V pk into $\geq 100\,\Omega$ diff.\ + transformer loss \\
	TX & PA & 175 (2\,V) & 350 & PAE $\geq 10$\,\% at $+16$\,dBm \\
	TX & LO & 13 & 14.3 & as RX \\
	\textbf{TX} & \textbf{total} & & \textbf{398} & BUD-09 $\leq 400$\,mW \\
	\bottomrule
\end{tabularx}\par\normalsize

\subsection{Budget and results}
\placeholder{parameter / budget / stub / design / status table}

% -----------------------------------------------------------------------------
\section{Changelog}

\begin{changelog}
	\change{2026-11-12}{integrator}{Requirements v0.1 (budget and interfaces) from design proposal}{Starting point for meeting~2}{all}{proposed}
	\change{2026-11-17}{integrator}{Requirements v0.1 $\to$ v0.2 (standard table)}{Changeover, see requirements\_changeover.md}{all}{proposed}
	\change{2026-11-27}{integrator}{Requirements v0.4: derived budget BUD-01\dots09, ADC/DAC models, interface levels, RX noise/linearity cascade, TX level plan, power budget}{Fill TBD values so that each block can design alone}{all}{proposed}
	\change{2026-11-27}{integrator}{Transformer centre tap: ``biases the receiver'' $\to$ secondary DC-floating, receiver biases its own gates}{No bias voltage may cross a block boundary (\specref{4})}{LNA, DMIX, UMIX, DRV, PA}{proposed}
	\change{2026-11-27}{integrator}{I8: transformer or C + R $\to$ transformer as UMIX load (= BPF tank)}{Rule~1: UMIX already has an inductive load}{UMIX, DRV}{proposed}
	\change{2026-11-27}{integrator}{Supply current and power: TBD $\to$ budget values (D10 reversed)}{Power budget is a team deliverable (\specref{3}, item~5)}{all}{proposed}
\end{changelog}

% =============================================================================

\chapter{System Integration}
\label{ch:integrator}

\blockowner{Wojciech Miśkowicz}

% -----------------------------------------------------------------------------
\section{Schematic}
Block diagram of the RX and TX chains with the shared LO. Every block-to-block
connection is labelled with its interface ID I1--I10 (table in
section~\ref{sec:integrator-interfaces}).

\figplaceholder{TRX block diagram with interfaces I1--I10}{fig:integrator-block-diagram}

% -----------------------------------------------------------------------------
\section{Requirements}
Source: project specification and requirements changeover v0.2.

\begin{requirements}
	\req{SYS-01}{Technology}{TSMC 40\,nm, \texttt{tsmcN40}, Cadence + SpectreRF}{\specref{2--3}}
	\req{SYS-02}{Operating frequency}{6\,GHz, 20\,MHz channel}{\specref{2--3}}
	\req{SYS-03}{Architecture}{direct conversion, one LO shared by RX and TX (TDD)}{\specref{2--3}}
	\req{SYS-04}{Supply}{VDD = 1.1\,V for all blocks; VDD\_PA = 2\,V}{\specref{2--3}}
	\req{SYS-05}{RX input / TX output}{50\,$\Omega$}{\specref{2--3}}
	\req{SYS-06}{RX noise figure}{$< 10$\,dB}{\specref{2--3}}
	\req{SYS-07}{RX input signal range}{$-82$\,dBm (sensitivity)\dots$-30$\,dBm}{\specref{2--3}}
	\req{SYS-08}{TX output power}{$+16$\,dBm into 50\,$\Omega$, 8\,dB back-off from $P_\mathrm{1dB}$}{\specref{2--3}}
	\req{SYS-09}{Test signals}{one tone (gain, noise), two tones (IP3), QPSK/QAM (LNA and PA required)}{\specref{2--3}}
	\req{SYS-10}{Devices and bias}{real \texttt{tsmcN40} devices, bias inside each block; baluns/transformers may be ideal with a realistic $Q$}{\specref{2--3}}
	\req{SYS-11}{Coupling rule}{see the coupling rules below and the interface table}{\specref{2--3}}
\end{requirements}

\subsection{Derived system budget}\label{sec:integrator-budget-req}
These rows are not in the specification. They follow from SYS-06\dots SYS-08 and
from the converter models in section~\ref{sec:integrator-converters}. Every block
requirement table refers back to them. The derivations are in
section~\ref{sec:integrator-budget}.

\begin{requirements}
	\req{BUD-01}{RX noise figure (budget)}{NF$_\mathrm{SSB} \leq 7.5$\,dB at maximum gain (estimate 7.3\,dB, 2.7\,dB margin to SYS-06); measured SSB with one tone, single (I) mixer}{SYS-06}
	\req{BUD-02}{RX linearity}{$P_\mathrm{1dB,in} \geq -22$\,dBm and IIP3 $\geq -12$\,dBm at minimum gain, referred to the RX input (estimate $-11.4$\,dBm)}{SYS-07}
	\req{BUD-03}{RX voltage gain (pin $\to$ ADC)}{21\dots77\,dB (LNA 15 + DMIX 8 + VGA $-2$\dots54); nominal AGC target 100\,mVrms diff.\ at the ADC}{SYS-07}
	\req{BUD-04}{RX IIP2}{$\geq +35$\,dBm referred to the RX input}{SYS-03}
	\req{BUD-05}{TX chain $OP_\mathrm{1dB}$}{$\geq +24$\,dBm at the antenna; every stage before the PA compresses $\geq 2$\,dB later than the PA (referred to the PA output)}{SYS-08}
	\req{BUD-06}{TX nominal level plan}{DAC 70\,mVrms diff.\ $\to$ $+16$\,dBm at the antenna (section~\ref{tab:integrator-tx-levels})}{SYS-08}
	\req{BUD-07}{LO amplitude at each mixer}{1.0\,Vpp diff.\ nominal, 0.8\dots1.2\,Vpp diff.\ over PVT}{SYS-03}
	\req{BUD-08}{Power, RX mode}{$\leq 45$\,mW from 1.1\,V (LNA + DMIX + VGA + LO)}{\specref{3}}
	\req{BUD-09}{Power, TX mode}{$\leq 400$\,mW (UMIX + DRV + LO from 1.1\,V, PA from 2\,V)}{\specref{3}}
\end{requirements}

\subsection{Coupling rules}
\begin{enumerate}
	\item If the driving stage already has an inductive load, that load is made a
	      \textbf{transformer/balun}. It also gives AC coupling and single-ended to
	      differential conversion. The \textbf{driving} block owns it.
	\item \textbf{Bias across a transformer.} The secondary of an interface transformer
	      is \textbf{DC-floating}: its centre tap is only AC-grounded (MIM capacitor
	      $\geq 2$\,pF, owned by the driving block) and has no DC path to VDD or ground.
	      The \textbf{receiving} block sets the DC level of its own input gates through
	      its own bias resistors ($\geq 5$\,k$\Omega$ per side, so they do not load the
	      tank). No bias voltage crosses a block boundary (\specref{4}).
	\item Otherwise use a \textbf{MIM capacitor plus bias resistor} in the \textbf{receiving} block.
	\item No coil is added only for coupling.
	\item Baseband: DMIX$\to$VGA is AC-coupled (removes the DC offset from LO
	      self-mixing). The baseband edges (VGA$\to$ADC, DAC$\to$UMIX) are
	      DC-coupled with a common mode of \textbf{0.55\,V}.
	\item Impedances: 50\,$\Omega$ ports are given as $Z$ plus $|S_{11}|$/$|S_{22}|$;
	      on-chip ports as $R \parallel C$. The \textbf{receiving} block owns the number;
	      the driving block meets its own specification into that load.
\end{enumerate}

\subsection{Interfaces}\label{sec:integrator-interfaces}
The agreed numbers at each interface are the contract between blocks. The LO drives
I3 and I7 together, so the LO output load is the sum of both mixer inputs.

\par\medskip\noindent\small
\begin{tabularx}{\textwidth}{@{}l l l l L l@{}}
	\toprule
	\textbf{ID} & \textbf{From $\to$ To} & \textbf{Freq.} & \textbf{Type} & \textbf{Coupling} & \textbf{Owner} \\
	\midrule
	I1 & Antenna (50\,$\Omega$) $\to$ LNA & 6\,GHz & SE & AC: series C in the LNA input match & LNA \\
	I2 & LNA $\to$ DMIX & 6\,GHz & SE $\to$ diff. & AC: \textbf{transformer used as the LNA load}; secondary DC-floating, DMIX biases its own gates & LNA \\
	I3 & LO $\to$ DMIX & 6\,GHz & diff. & AC: C + R at the DMIX LO input & DMIX \\
	I4 & DMIX $\to$ VGA & BB $\leq 10$\,MHz & diff. & AC: C + R, $f_\mathrm{c} \leq 100$\,kHz & VGA \\
	I5 & VGA $\to$ ADC load & BB & diff. & \textbf{DC}, common mode 0.55\,V & -- \\
	I6 & DAC $\to$ UMIX & BB & diff. & \textbf{DC}, common mode 0.55\,V & -- \\
	I7 & LO $\to$ UMIX & 6\,GHz & diff. & AC: C + R at the UMIX LO input & UMIX \\
	I8 & UMIX $\to$ DRV & 6\,GHz & diff. & AC: \textbf{transformer used as the UMIX load} (= BPF tank, rule~1); secondary DC-floating. Proposed, UMIX and DRV owners confirm & UMIX \\
	I9 & DRV $\to$ PA & 6\,GHz & diff. & AC: \textbf{transformer used as the DRV load} (resonates the large PA $C_\mathrm{in}$); secondary DC-floating, PA biases its own gates & DRV \\
	I10 & PA $\to$ antenna (50\,$\Omega$) & 6\,GHz & diff. $\to$ SE & output transformer / matching network & PA \\
	\bottomrule
\end{tabularx}\par\normalsize

\medskip
\noindent Levels at each interface (CW test tone, nominal; ``diff.'' = differential,
``pk'' = peak, ``Vpp'' = peak-to-peak):

\par\medskip\noindent\small
\begin{tabularx}{\textwidth}{@{}l P{0.17\textwidth} P{0.25\textwidth} L@{}}
	\toprule
	\textbf{ID} & \textbf{DC level} & \textbf{Signal (min \dots max)} & \textbf{Load (owned by the receiver)} \\
	\midrule
	I1 & LNA gate $\approx 0.45$\,V (internal) & $-82$\dots$-30$\,dBm (50\,$\mu$Vpp\dots20\,mVpp at the pin) & 50\,$\Omega$, $|S_{11}| \leq -10$\,dB \\
	I2 & 0.55\,V, set by DMIX & 0.28\dots112\,mVpp diff. & $\geq 2$\,k$\Omega$ $\parallel$ $\leq 80$\,fF diff. \\
	I3, I7 & set by each mixer & 1.0\,Vpp diff.\ (0.8\dots1.2) & per mixer $\geq 5$\,k$\Omega$ $\parallel$ $\leq 80$\,fF diff. \\
	I4 & DMIX $\approx 0.7$\,V; VGA 0.55\,V & 0.71\dots283\,mVpp diff. & $\geq 50$\,k$\Omega$ $\parallel$ $\leq 1$\,pF diff.\ (after 20\,pF series C per side) \\
	I5 & \textbf{0.55\,V $\pm$ 30\,mV} & 283\,mVpp diff.\ (100\,mVrms) nominal, 1.0\,Vpp diff.\ full scale & ADC: 10\,k$\Omega$ diff.\ $\parallel$ 1\,pF per side \\
	I6 & \textbf{0.55\,V $\pm$ 50\,mV} & 200\,mVpp diff.\ (70\,mVrms) nominal, 1.0\,Vpp diff.\ full scale & $\geq 10$\,k$\Omega$ $\parallel$ $\leq 1$\,pF diff.; DAC $R_\mathrm{s} = 200\,\Omega$ diff. \\
	I8 & DRV gate $\approx 0.5$\,V, set by DRV & 0.2\,Vpp diff.\ (0.1\,V pk) & $\geq 1$\,k$\Omega$ $\parallel$ $\leq 150$\,fF diff. \\
	I9 & PA gate $\approx 0.45$\,V, set by PA & 1.0\,Vpp diff.\ (0.5\,V pk) & $\geq 100\,\Omega$ $\parallel$ $\leq 1.0$\,pF diff. \\
	I10 & drain at VDD\_PA = 2\,V & $+16$\,dBm (2.0\,V pk); $OP_\mathrm{1dB} \geq +24$\,dBm & 50\,$\Omega$, $|S_{22}| \leq -10$\,dB \\
	\bottomrule
\end{tabularx}\par\normalsize

% -----------------------------------------------------------------------------
\section{Overview}

\subsection{Top-level testbench}
\placeholder{\texttt{trx\_top} hierarchy, config views, stimulus and loads}

\figplaceholder{Top-level testbench \texttt{trx\_top}}{fig:integrator-testbench}

\subsection{Converter models at the baseband edges}\label{sec:integrator-converters}
The project has no ADC or DAC. The testbench models them as below. The values
are typical for a 20\,MHz Wi-Fi transceiver in a 1.1\,V, 40\,nm process.

\par\medskip\noindent\small
\begin{tabularx}{\textwidth}{@{}P{0.3\textwidth} L L@{}}
	\toprule
	\textbf{Parameter} & \textbf{ADC (RX, I5)} & \textbf{DAC (TX, I6)} \\
	\midrule
	Type & 10-bit SAR, 80\,MS/s (2$\times$ oversampled) & 10-bit current-steering, 160\,MS/s, with reconstruction filter \\
	Full scale & 1.0\,Vpp diff.\ (0.5\,V pk diff.) & 1.0\,Vpp diff. \\
	Common mode & 0.55\,V $\pm$ 30\,mV accepted & 0.55\,V $\pm$ 50\,mV delivered \\
	ENOB / SNR & 9 bit / 56\,dB $\to$ 89\,nV/$\sqrt{\mathrm{Hz}}$ diff.\ input-referred & 9 bit \\
	Nominal level & 100\,mVrms diff.\ ($-11$\,dBFS): OFDM PAPR 10\,dB + 4\,dB AGC margin & 70\,mVrms diff.\ ($-14$\,dBFS): set by UMIX linearity \\
	Model in \texttt{trx\_top} & 5\,k$\Omega$ per side to 0.55\,V + 1\,pF per side to ground ($R \approx 1/(f_\mathrm{s} C_\mathrm{s})$) & ideal source, 100\,$\Omega$ per side, CM 0.55\,V \\
	\bottomrule
\end{tabularx}\par\normalsize

\subsection{RX budget}\label{sec:integrator-budget}

\paragraph{Gain and level plan.}
Voltage gains are defined at the pins: LNA from the matched input pin (single-ended)
to its differential output; mixer conversion gain from RF diff.\ amplitude to BB diff.\
amplitude. The AGC target at the ADC is 100\,mVrms diff.\ (CW-equivalent).

\par\medskip\noindent\small
\begin{tabularx}{\textwidth}{@{}l c L L L@{}}
	\toprule
	\textbf{Node} & \textbf{Gain} & \textbf{$-82$\,dBm} & \textbf{$-30$\,dBm} & \textbf{$-22$\,dBm ($P_\mathrm{1dB}$ target)} \\
	\midrule
	RX input pin (SE) & -- & 50\,$\mu$Vpp & 20\,mVpp & 50\,mVpp \\
	LNA out / DMIX in (diff.) & 15\,dB & 0.28\,mVpp & 112\,mVpp & 283\,mVpp \\
	DMIX out / VGA in (diff.) & 8\,dB & 0.71\,mVpp & 283\,mVpp & 710\,mVpp \\
	VGA out / ADC (diff.) & $+52$ / $0$ / $0$\,dB & 283\,mVpp & 283\,mVpp & 710\,mVpp (FS 1.0\,Vpp) \\
	\bottomrule
\end{tabularx}\par\normalsize

\medskip
\noindent The required VGA range is 52\,dB (0\dots52\,dB). The VGA gives
$-2$\dots$+54$\,dB, so 2\,dB margin is left at each end for the PVT spread of the
LNA and DMIX gain. The LNA has one gain mode (decision D5), so the whole range is in
the VGA. At $-22$\,dBm with the gain still set for $-30$\,dBm the signal is 3\,dB
below the ADC full scale, so the ADC does not clip before the analog chain compresses.

\paragraph{Noise.}
All contributions are referred to the RX input pin as noise voltage density.
The source term is $kTR_\mathrm{s} = 2.0\cdot10^{-19}$\,V$^2$/Hz (0.447\,nV/$\sqrt{\mathrm{Hz}}$,
matched 50\,$\Omega$, $T_0 = 290$\,K). A block with input-referred density $v_n$
after voltage gain $G$ contributes $v_n^2/G^2$. NF is SSB (one mixer, the noise of
both sidebands folds onto the wanted one), which is 3\,dB worse than the DSB value
that an I/Q receiver would show. VGA at maximum gain:

\par\medskip\noindent\small
\begin{tabularx}{\textwidth}{@{}l l l L r@{}}
	\toprule
	\textbf{Block} & \textbf{Gain before} & \textbf{Own noise (budget)} & \textbf{Referred to input pin} & \textbf{Share} \\
	\midrule
	Source 50\,$\Omega$ & -- & -- & 0.447\,nV/$\sqrt{\mathrm{Hz}}$ & 19\,\% \\
	LNA & 0\,dB & NF $\leq 2$\,dB & 0.342\,nV/$\sqrt{\mathrm{Hz}}$ & 11\,\% \\
	DMIX & 15\,dB & $\leq 4.0$\,nV/$\sqrt{\mathrm{Hz}}$ diff.\ (SSB) & 0.711\,nV/$\sqrt{\mathrm{Hz}}$ & 47\,\% \\
	VGA + LPF & 23\,dB & $\leq 7.0$\,nV/$\sqrt{\mathrm{Hz}}$ diff. & 0.496\,nV/$\sqrt{\mathrm{Hz}}$ & 23\,\% \\
	ADC & 77\,dB & 89\,nV/$\sqrt{\mathrm{Hz}}$ diff. & 0.013\,nV/$\sqrt{\mathrm{Hz}}$ & $< 0.1$\,\% \\
	\midrule
	\textbf{Chain} & & & \textbf{NF$_\mathrm{SSB}$ = 7.3\,dB} (DSB $\approx 4.3$\,dB) & 100\,\% \\
	\bottomrule
\end{tabularx}\par\normalsize

\medskip
\noindent The input noise floor in 20\,MHz is $-93.7$\,dBm, so the SNR at $-82$\,dBm is
11.7\,dB. With a single-mode 15\,dB LNA the DMIX dominates (47\,\%), so its
noise budget is the tightest one in the RX. The 2.7\,dB margin to SYS-06 covers the
transformer loss at I2, real inductor $Q$ in the LNA and the PVT corners. At minimum VGA
gain the NF rises to $\approx 28$\,dB, but the input is then $\geq -30$\,dBm and the
SNR is still $\approx 43$\,dB.

\paragraph{Linearity.}
Two-tone IIP3 of each block is referred to the RX input and summed as
$1/\mathrm{IIP3} = \sum 1/\mathrm{IIP3}_{i,\mathrm{ref}}$ (in mW). The worst case is
the minimum VGA gain, where the full $-30$\,dBm signal reaches the VGA. Block IIP3 on
on-chip nodes is given as the \textbf{per-tone peak differential voltage at the block
input} at the intercept point.

\par\medskip\noindent\small
\begin{tabularx}{\textwidth}{@{}l L l r@{}}
	\toprule
	\textbf{Block} & \textbf{Requirement} & \textbf{Referred to RX input} & \textbf{Share} \\
	\midrule
	LNA & IIP3 $\geq -4$\,dBm (50\,$\Omega$) & $-4.0$\,dBm & 18\,\% \\
	DMIX & IIP3 $\geq 0.7$\,V pk diff.\ per tone at the RF input & $-8.1$\,dBm & 46\,\% \\
	VGA (code 0, $-2$\,dB) & OIP3 $\geq 1.6$\,V pk diff.\ per tone & $-6.9$\,dBm & 35\,\% \\
	\midrule
	\textbf{Chain} & & \textbf{$-11.4$\,dBm} & 100\,\% \\
	\bottomrule
\end{tabularx}\par\normalsize

\medskip
\noindent The requirement is IIP3 $\geq -12$\,dBm, from $P_\mathrm{1dB,in} \geq -22$\,dBm
($-30$\,dBm + 8\,dB) and IIP3 $\approx P_\mathrm{1dB} + 9.6$\,dB. The margin is only
0.6\,dB, so \textbf{linearity, not noise, is the critical RX budget}. If the DMIX
cannot reach 0.7\,V pk, the first remedy is a switchable LNA gain step
(e.g.\ 15/3\,dB, revisiting D5). It costs no noise at sensitivity and gains
$\approx 6$\,dB of IIP3.

\paragraph{IIP2 and DC offset.}
In direct conversion an AM interferer gives IM2 at DC and in band. BUD-04
($+35$\,dBm) keeps IM2 from a $-30$\,dBm two-tone signal 65\,dB below the signal. In
practice the DMIX mismatch sets IIP2. The DC offset from LO self-mixing is blocked
by the AC coupling at I4 ($f_\mathrm{c} \leq 100$\,kHz, below the first Wi-Fi
subcarrier at 312.5\,kHz). The VGA must cancel its own offset, because at 54\,dB a
2\,mV input offset would saturate the output.

\paragraph{Baseband filter.}
The LPF is a 3rd-order Butterworth with $f_{-3\mathrm{dB}} = 12$\,MHz. It costs
0.9\,dB at 9.4\,MHz (edge of the occupied band), gives 13.5\,dB at 20\,MHz (adjacent
channel) and 46\,dB at 70\,MHz (the first ADC alias at $f_\mathrm{s} - 10$\,MHz).

\subsection{TX budget}\label{tab:integrator-tx-levels}
The TX is specified in voltages up to the PA input and in dBm at the antenna. Each
stage before the PA has $\geq 10$\,dB back-off at the nominal level, so the PA alone
sets the chain $P_\mathrm{1dB}$. Without I/Q both sidebands ($f_\mathrm{LO} \pm f_\mathrm{BB}$)
leave the mixer at equal level. The $+16$\,dBm applies to the wanted sideband, and the
PA and DRV are also characterised alone with an RF source.

\par\medskip\noindent\small
\begin{tabularx}{\textwidth}{@{}l l L L@{}}
	\toprule
	\textbf{Node} & \textbf{Gain} & \textbf{Nominal (CW)} & \textbf{Compression requirement} \\
	\midrule
	DAC out / UMIX in & -- & 70\,mVrms = 0.1\,V pk diff. & FS 0.5\,V pk diff. \\
	UMIX + BPF out / DRV in & 0\,dB (wanted sideband) & 0.1\,V pk diff. & $OP_\mathrm{1dB} \geq 0.4$\,V pk diff.\ (12\,dB back-off) \\
	DRV out / PA in & 14\,dB & 0.5\,V pk diff. & $OP_\mathrm{1dB} \geq 1.6$\,V pk diff.\ (10\,dB back-off) \\
	PA out (50\,$\Omega$, SE) & 12\,dB (2.0\,V/0.5\,V) & $+16$\,dBm = 2.0\,V pk & $OP_\mathrm{1dB} \geq +24$\,dBm (5.0\,V pk); target $+25$\,dBm \\
	\bottomrule
\end{tabularx}\par\normalsize

\medskip
\noindent At the PA $P_\mathrm{1dB}$ its input is $0.5 \cdot 10^{9/20} = 1.41$\,V pk diff.
The DRV must deliver this with $\geq 1$\,dB to spare, so it needs 1.6\,V pk. The DAC sits
at $-14$\,dBFS. A 10-bit DAC still gives $\approx 42$\,dB SNR there, well above the
$\approx 25$\,dB needed for 64-QAM, and the UMIX input stays small enough for the
0.55\,V common mode at 1.1\,V supply.

\subsection{Power budget}

\par\medskip\noindent\small
\begin{tabularx}{\textwidth}{@{}l l r r L@{}}
	\toprule
	\textbf{Mode} & \textbf{Block} & \textbf{$I$ [mA]} & \textbf{$P$ [mW]} & \textbf{Basis} \\
	\midrule
	RX & LNA & 8 & 8.8 & NF 2\,dB at 6\,GHz needs $g_m \approx 50$--60\,mS \\
	RX & DMIX & 6 & 6.6 & Gilbert cell, $\approx 2.5$\,mA per side in the $g_m$ stage \\
	RX & VGA + LPF & 10 & 11 & 3 stages + LPF + offset cancellation \\
	RX & LO & 13 & 14.3 & LO-16/17 (owner's value) \\
	\textbf{RX} & \textbf{total} & \textbf{37} & \textbf{40.7} & BUD-08 $\leq 45$\,mW \\
	\midrule
	TX & UMIX & 6 & 6.6 & as DMIX \\
	TX & DRV & 25 & 27.5 & class A, 1.6\,V pk into $\geq 100\,\Omega$ diff.\ + transformer loss \\
	TX & PA & 175 (2\,V) & 350 & PAE $\geq 10$\,\% at $+16$\,dBm \\
	TX & LO & 13 & 14.3 & as RX \\
	\textbf{TX} & \textbf{total} & & \textbf{398} & BUD-09 $\leq 400$\,mW \\
	\bottomrule
\end{tabularx}\par\normalsize

\subsection{Budget and results}
\placeholder{parameter / budget / stub / design / status table}

% -----------------------------------------------------------------------------
\section{Changelog}

\begin{changelog}
	\change{2026-11-12}{integrator}{Requirements v0.1 (budget and interfaces) from design proposal}{Starting point for meeting~2}{all}{proposed}
	\change{2026-11-17}{integrator}{Requirements v0.1 $\to$ v0.2 (standard table)}{Changeover, see requirements\_changeover.md}{all}{proposed}
	\change{2026-11-27}{integrator}{Requirements v0.4: derived budget BUD-01\dots09, ADC/DAC models, interface levels, RX noise/linearity cascade, TX level plan, power budget}{Fill TBD values so that each block can design alone}{all}{proposed}
	\change{2026-11-27}{integrator}{Transformer centre tap: ``biases the receiver'' $\to$ secondary DC-floating, receiver biases its own gates}{No bias voltage may cross a block boundary (\specref{4})}{LNA, DMIX, UMIX, DRV, PA}{proposed}
	\change{2026-11-27}{integrator}{I8: transformer or C + R $\to$ transformer as UMIX load (= BPF tank)}{Rule~1: UMIX already has an inductive load}{UMIX, DRV}{proposed}
	\change{2026-11-27}{integrator}{Supply current and power: TBD $\to$ budget values (D10 reversed)}{Power budget is a team deliverable (\specref{3}, item~5)}{all}{proposed}
\end{changelog}
